\section{Proof of the trace-support theorem}\label{sec:conclusion}

We now sum the curvature estimate against the cyclic tuple weights.
The order of the two limiting choices is essential: the spatial radius
is fixed first, uniformly in the conjugacy class, and only then does
the simplex dimension tend to infinity.

\begin{proof}[Proof of Theorem~\ref{thm:main}]
By Proposition~\ref{red:reduction}, it suffices to treat a finitely
generated group with a fixed word metric and an idempotent $p$ having
an exponential moment. The zero idempotent is immediate, so assume
$M_p>0$, with $M_p$ as in Lemma~\ref{cy:weights-lemma}.

For the construction of Sections~\ref{sec:connection}
and~\ref{sec:curvature}, write
\[
 Q_R=C S^{3/2}R^{-1/10},\qquad S=b_*+2.
\]
The fixed constants here may depend on the word metric, but not on
$g$, the tuple, or $n$. Since $S=O(1+\log R)$, we have $Q_R\to0$.
Choose $R$ once and for all so large that the construction applies and
\begin{equation}\label{final:radius-choice}
 Q_RM_p<1.
\end{equation}
This fixes the final spatial radius $r=6^{J_*-1}R$, the scale ladder,
and all constants denoted $C_R$.

Let $g$ have infinite order, or have finite order with
$\min_h d(h,gh)>4r$. Proposition~\ref{co:connection} provides a
compatible connection family through every sufficiently large even
dimension $n$. Combining Proposition~\ref{cy:trace-identity},
Theorem~\ref{pf:bound}, and \eqref{cy:moment}, we obtain
\begin{align}
 |\tau_{[g]}(p)|
 &\le \sum_{h\in\mathcal T_n(g)}|v_{n+1}(h)|
               \mathbb E_{\Delta_n}|\Pf(d\alpha_h)|\notag\\
 &\le C_RQ_R^n\sum_{h\in\mathcal T_n(g)}|v_{n+1}(h)|
                                    \prod_{i=1}^n l_i\notag\\
 &\le C_RM_p\,(Q_RM_p)^n.\label{final:decay}
\end{align}
In the last step we used $l_0\ge1$ to insert the missing closing-edge
factor. The constant $C_R$ is independent of $n$ and $g$.
Letting $n$ tend to infinity through even integers in
\eqref{final:decay} gives $\tau_{[g]}(p)=0$.

It remains to identify the possible surviving classes. Every
infinite-order class has just been eliminated. If a finite-order class
does not satisfy the displacement condition, some conjugate of its
representative has word length at most $4r$, because
\[
 d(h,gh)=|h^{-1}gh|.
\]
Thus each surviving class meets the finite ball $B(1,4r)$. There are
only finitely many such classes. This proves the result for $p$.
Proposition~\ref{red:reduction} transfers it to the original $e$ and
group. Finally, an element of $K_0(\ell^1(G))$ is a difference of
idempotent classes; taking the union of their two finite supports
proves the asserted inclusion for $\operatorname{HS}^1$.
\end{proof}
